\documentclass{article}
\usepackage[utf8]{inputenc}
\usepackage{booktabs}
\usepackage{longtable}
\title{Full local review smoke}
\begin{document}
\maketitle
\begin{abstract}
Full local review smoke synthesizes 3 papers across 6 evidence rows using factor structure as the main organizing lens. The synthesis is concentrated in design factors, llm-assisted-extraction, with recurring methods such as pipeline, llm-assisted-extraction and task coverage centered on literature review, information extraction. Across 3 papers, the evidence is most coherently organized around factor structure rather than any single method, and the strongest clusters are design factors, llm-assisted-extraction, pipeline. The main unresolved gap is that the current literature set shows weak metric reporting, limiting rigorous cross-paper evaluation., which keeps broader comparison and transfer claims provisional.
\end{abstract}
\noindent\textbf{Keywords:} information extraction, llm-assisted-extraction, pipeline, literature review, automatically extract structured literature review signals from academic papers., factor taxonomy, design factors, current set
\section{Introduction and Review Scope}
Introduction and Review Scope provides the analytic entry point for this subsection, bringing the discussion into focus around Design Factors, Llm Assisted Extraction. Across 3 evidence units, the relevant literature assembles method families such as Pipeline; Llm Assisted Extraction, tasks such as Literature Review, Information Extraction, which defines the comparison set for the discussion. Read this way, the subsection begins from a shared problem space rather than a sequence of isolated paper summaries. \cite{02_2003_brosseau_trf,01_2001_hamed_safety_science}

Across the reviewed studies, a recurring pattern emerges in introduction and review scope. The system supports downstream gap analysis by preserving evidence links. The system supports downstream gap analysis by preserving evidence links. Collectively, these observations support a subsection-level claim instead of leaving the evidence as a descriptive inventory. \cite{02_2003_brosseau_trf,01_2001_hamed_safety_science,03_2006_yang_tra}

Taken together, the literature on introduction and review scope points to a field-level pattern rather than a single-study result. The evidence clusters most clearly around Design Factors, Llm Assisted Extraction. The synthesis therefore clarifies both the main takeaway and the boundary conditions that qualify it. \cite{02_2003_brosseau_trf,01_2001_hamed_safety_science}

An unresolved problem in this literature concerns The current literature set shows weak metric reporting, limiting rigorous cross-paper evaluation. The available evidence points to the issue, but the reporting pattern remains too partial to sustain strong generalization. Until that gap is addressed, any synthesis of introduction and review scope remains conditional rather than fully generalizable. \cite{02_2003_brosseau_trf,01_2001_hamed_safety_science}

\section{Factor Taxonomy and Contextual Drivers}
Factor Taxonomy and Contextual Drivers provides the analytic entry point for this subsection, bringing the discussion into focus around Design Factors, Llm Assisted Extraction. Across 3 evidence units, the relevant literature assembles method families such as Pipeline; Llm Assisted Extraction, tasks such as Literature Review, Information Extraction, which defines the comparison set for the discussion. Read this way, the subsection begins from a shared problem space rather than a sequence of isolated paper summaries. \cite{02_2003_brosseau_trf,01_2001_hamed_safety_science}

Across the reviewed studies, a recurring pattern emerges in factor taxonomy and contextual drivers. The system supports downstream gap analysis by preserving evidence links. The system supports downstream gap analysis by preserving evidence links. Collectively, these observations support a subsection-level claim instead of leaving the evidence as a descriptive inventory. \cite{02_2003_brosseau_trf,01_2001_hamed_safety_science,03_2006_yang_tra}

Comparison is most informative where work on factor taxonomy and contextual drivers converges in broad outline but diverges in operational detail. design factors appears in 6 evidence unit(s), with 0 gap signal(s) and 0 contradiction signal(s). llm-assisted-extraction appears in 6 evidence unit(s), with 0 gap signal(s) and 0 contradiction signal(s). The resulting contrast shows where convergence is substantive and where apparent agreement rests on non-equivalent designs. \cite{02_2003_brosseau_trf,01_2001_hamed_safety_science,03_2006_yang_tra}

Taken together, the literature on factor taxonomy and contextual drivers points to a field-level pattern rather than a single-study result. The evidence clusters most clearly around Design Factors, Llm Assisted Extraction. The synthesis therefore clarifies both the main takeaway and the boundary conditions that qualify it. \cite{02_2003_brosseau_trf,01_2001_hamed_safety_science}

An unresolved problem in this literature concerns The current literature set shows weak metric reporting, limiting rigorous cross-paper evaluation. The available evidence points to the issue, but the reporting pattern remains too partial to sustain strong generalization. Until that gap is addressed, any synthesis of factor taxonomy and contextual drivers remains conditional rather than fully generalizable. \cite{02_2003_brosseau_trf,01_2001_hamed_safety_science}

\section{Factor Focus: Design Factors}
Design Factors provides the analytic entry point for this subsection, bringing the discussion into focus around Design Factors, Llm Assisted Extraction. Across 3 evidence units, the relevant literature assembles method families such as Pipeline; Llm Assisted Extraction, tasks such as Literature Review, Information Extraction, which defines the comparison set for the discussion. Read this way, the subsection begins from a shared problem space rather than a sequence of isolated paper summaries. \cite{02_2003_brosseau_trf,01_2001_hamed_safety_science}

Across the reviewed studies, a recurring pattern emerges in design factors. The system supports downstream gap analysis by preserving evidence links. The system supports downstream gap analysis by preserving evidence links. Collectively, these observations support a subsection-level claim instead of leaving the evidence as a descriptive inventory. \cite{02_2003_brosseau_trf,01_2001_hamed_safety_science,03_2006_yang_tra}

Taken together, the literature on design factors points to a field-level pattern rather than a single-study result. The evidence clusters most clearly around Design Factors, Llm Assisted Extraction. The synthesis therefore clarifies both the main takeaway and the boundary conditions that qualify it. \cite{02_2003_brosseau_trf,01_2001_hamed_safety_science}

An unresolved problem in this literature concerns The current literature set shows weak metric reporting, limiting rigorous cross-paper evaluation. The available evidence points to the issue, but the reporting pattern remains too partial to sustain strong generalization. Until that gap is addressed, any synthesis of design factors remains conditional rather than fully generalizable. \cite{02_2003_brosseau_trf,01_2001_hamed_safety_science}

\section{Comparative Factor Evidence and Interactions}
Comparative Factor Evidence and Interactions provides the analytic entry point for this subsection, bringing the discussion into focus around Design Factors, Llm Assisted Extraction. Across 3 evidence units, the relevant literature assembles method families such as Pipeline; Llm Assisted Extraction, tasks such as Literature Review, Information Extraction, which defines the comparison set for the discussion. Read this way, the subsection begins from a shared problem space rather than a sequence of isolated paper summaries. \cite{02_2003_brosseau_trf,01_2001_hamed_safety_science}

Comparison is most informative where work on comparative factor evidence and interactions converges in broad outline but diverges in operational detail. design factors appears in 6 evidence unit(s), with 0 gap signal(s) and 0 contradiction signal(s). llm-assisted-extraction appears in 6 evidence unit(s), with 0 gap signal(s) and 0 contradiction signal(s). The resulting contrast shows where convergence is substantive and where apparent agreement rests on non-equivalent designs. \cite{02_2003_brosseau_trf,01_2001_hamed_safety_science,03_2006_yang_tra}

Across the reviewed studies, a recurring pattern emerges in comparative factor evidence and interactions. The system supports downstream gap analysis by preserving evidence links. The system supports downstream gap analysis by preserving evidence links. Collectively, these observations support a subsection-level claim instead of leaving the evidence as a descriptive inventory. \cite{02_2003_brosseau_trf,01_2001_hamed_safety_science,03_2006_yang_tra}

Taken together, the literature on comparative factor evidence and interactions points to a field-level pattern rather than a single-study result. The evidence clusters most clearly around Design Factors, Llm Assisted Extraction. The synthesis therefore clarifies both the main takeaway and the boundary conditions that qualify it. \cite{02_2003_brosseau_trf,01_2001_hamed_safety_science}

An unresolved problem in this literature concerns The current literature set shows weak metric reporting, limiting rigorous cross-paper evaluation. The available evidence points to the issue, but the reporting pattern remains too partial to sustain strong generalization. Until that gap is addressed, any synthesis of comparative factor evidence and interactions remains conditional rather than fully generalizable. \cite{02_2003_brosseau_trf,01_2001_hamed_safety_science}

\section{Research Gaps and Opportunities}
Research Gaps and Opportunities provides the analytic entry point for this subsection, bringing the discussion into focus around Design Factors, Llm Assisted Extraction. Across 3 evidence units, the relevant literature assembles method families such as Pipeline; Llm Assisted Extraction, tasks such as Literature Review, Information Extraction, which defines the comparison set for the discussion. Read this way, the subsection begins from a shared problem space rather than a sequence of isolated paper summaries. \cite{02_2003_brosseau_trf,01_2001_hamed_safety_science}

An unresolved problem in this literature concerns The current literature set shows weak metric reporting, limiting rigorous cross-paper evaluation. The available evidence points to the issue, but the reporting pattern remains too partial to sustain strong generalization. Until that gap is addressed, any synthesis of research gaps and opportunities remains conditional rather than fully generalizable. \cite{02_2003_brosseau_trf,01_2001_hamed_safety_science}

Taken together, the literature on research gaps and opportunities points to a field-level pattern rather than a single-study result. The evidence clusters most clearly around Design Factors, Llm Assisted Extraction. The synthesis therefore clarifies both the main takeaway and the boundary conditions that qualify it. \cite{02_2003_brosseau_trf,01_2001_hamed_safety_science}

\section{Conclusion}
The field-level picture across 3 papers is that factor structure provide the clearest organizing lens, with the strongest signals concentrated in design factors, llm-assisted-extraction, pipeline. Stable conclusions: Across 3 papers, the evidence is most coherently organized around factor structure rather than any single method, and the strongest clusters are design factors, llm-assisted-extraction, pipeline. Methodological diversity is real, but the review still shows recurring families such as pipeline, llm-assisted-extraction, which suggests the field is comparing similar problems with different operational choices. Reporting remains uneven across tasks and metrics such as literature review, information extraction, which limits direct comparison across studies. Unresolved tensions: One unresolved tension is that The current literature set shows weak metric reporting, limiting rigorous cross-paper evaluation.. A second tension is that current studies still need common evaluation protocols and explicit metric reporting conventions that let future reviews compare outcomes on the same scale. Research priorities: Priority should be given to common evaluation protocols and explicit metric reporting conventions that let future reviews compare outcomes on the same scale. Future work should directly target The current literature set shows weak metric reporting, limiting rigorous cross-paper evaluation.. A third priority is to align methods, metrics, and context descriptors so the literature can be synthesized more defensibly. \cite{02_2003_brosseau_trf}

\appendix
\section{Appendix}
\subsection{Appendix Summary}
\noindent The appendix records 3 papers and 6 matrix rows used to generate the review synthesis. The most visible methodological families are pipeline, llm-assisted-extraction, while recurring tasks include literature review, information extraction. Verified gaps are indexed separately so readers can trace which evidence clusters remain unresolved before moving to the main conclusion.

\begin{itemize}
\item \textbf{Papers}: 3
\item \textbf{Matrix rows}: 6
\item \textbf{Verified gaps}: 1
\item \textbf{Dominant axis}: factor
\item \textbf{Top methods}: pipeline, llm-assisted-extraction
\item \textbf{Top tasks}: literature review, information extraction
\end{itemize}

\subsection{Evidence Table}
\begin{longtable}{p{0.14\textwidth}p{0.21\textwidth}p{0.10\textwidth}p{0.17\textwidth}p{0.17\textwidth}p{0.17\textwidth}}
\toprule
Paper & Title & Year & Methods & Tasks & Gap matches \\
\midrule
\endfirsthead
\toprule
Paper & Title & Year & Methods & Tasks & Gap matches \\
\midrule
\endhead
03\_2006\_yang\_tra & Modeling pedestrians’ road crossing behavior in traffic system & 2006 & pipeline, llm-assisted-extraction & literature review, information extraction & e3ceb20c-9aef-4214-b4d6-e84da48999fb \\
02\_2003\_brosseau\_trf & crosswalks and a pedestrian user survey which sought pedestrian perceptions toward various pedestrian & 2003 & pipeline, llm-assisted-extraction & literature review, information extraction & e3ceb20c-9aef-4214-b4d6-e84da48999fb \\
01\_2001\_hamed\_safety\_science & Untitled & 2001 & pipeline, llm-assisted-extraction & literature review, information extraction & e3ceb20c-9aef-4214-b4d6-e84da48999fb \\
\bottomrule
\end{longtable}

\subsection{Gap Index}
\begin{longtable}{p{0.12\textwidth}p{0.44\textwidth}p{0.14\textwidth}p{0.22\textwidth}}
\toprule
Gap & Statement & Severity & Research need \\
\midrule
\endfirsthead
\toprule
Gap & Statement & Severity & Research need \\
\midrule
\endhead
e3ceb20c-9aef-4214-b4d6-e84da48999fb & The current literature set shows weak metric reporting, limiting rigorous cross-paper evaluation. & -- & common evaluation protocols and explicit metric reporting conventions that let future reviews compare outcomes on the same scale \\
\bottomrule
\end{longtable}
\end{document}
% BIB START
@article{01_2001_hamed_safety_science,
  author={Hamed},
  title={The pipeline can transform parsed paper chunks into a structured paper profile.},
  year={2001},
  journal={safety science}
}

@article{02_2003_brosseau_trf,
  author={Brosseau},
  title={crosswalks and a pedestrian user survey which sought pedestrian perceptions toward various pedestrian},
  year={2003},
  journal={trf}
}

@article{03_2006_yang_tra,
  author={Yang},
  title={Modeling pedestrians’ road crossing behavior in traffic system},
  year={2006},
  journal={tra}
}